% GENERATED FILE. Edit canonical lessons, then run npm run generate:books.
% canonical-source-hash: fnv1a64:90eb61feb22bb19f
% canonical-lessons: FR-C283-diffuser, FR-C283-diffusion, FR-C283-filmer, FR-C283-interviewer, FR-C283-antenne, FR-R283-revision-276-283

\chapter{À l'antenne}
\label{ch:fr-a-l-antenne}
\hlchaptermodality{283}

\begin{chapteropening}
\textbf{By the end of this chapter:} \emph{I can say that something was broadcast, filmed or interviewed, and that someone spoke on air.}

Complete the last lesson of chapter 283: I can say that something was broadcast, filmed or interviewed, and that someone spoke on air.
\end{chapteropening}

\section[diffuser]{diffuser — to broadcast; to spread}

\label{lesson:FR-C283-diffuser}



\begin{quote}
Before the new one: say the French for the editorial, then the French for short news items.
\end{quote}

\subsection*{You'll want to know: diffuser}
\textbf{diffuser} — \textquotedblleft{}to broadcast; to spread\textquotedblright{}.

\begin{quote}
La télé a diffusé le match en direct.
\end{quote}

\begin{quote}
\emph{The TV broadcast the match live.}
\end{quote}

Also to spread news or a smell: \emph{diffuser une information}.

\begin{cousinweb}
\emph{diffuser} is built on \emph{diffus} \textquotedblleft{}spread out\textquotedblright{} (EtymDB), a learned borrowing from Latin \emph{diffūsus}, the past participle of \emph{diffundere} \textquotedblleft{}to pour out\textquotedblright{} (Lewis \& Short).
\end{cousinweb}

\begin{grammarlens}[title={What you've built}]
One more word for telling, explaining and describing at B1.
\end{grammarlens}

\subsection*{Your turn}
\begin{itemize}
\raggedright
  \item \emph{Say it:} \emph{diffuser}
  \item \emph{Say it:} \emph{diffuser}, once more
  \item \emph{Say it:} a sentence of your own with \emph{diffuser}
  \item \emph{Recall:} say the French for the editorial, then the French for short news items, then say \emph{diffuser} again
\end{itemize}

\begin{tcolorbox}[breakable,colback=teal!4,colframe=teal!35!black,title={Before you move on}]
What does diffuser mean? (\textquotedblleft{}To broadcast; to spread\textquotedblright{}.) Say it once more.
\end{tcolorbox}
\section[la diffusion]{la diffusion — the broadcasting, the showing; the spread}

\label{lesson:FR-C283-diffusion}



\begin{quote}
Before the new one: say the French for short news items, then the French for to broadcast.
\end{quote}

\subsection*{You'll want to know: la diffusion}
\textbf{la diffusion} — \textquotedblleft{}the broadcasting, the showing; the spread\textquotedblright{}.

\begin{quote}
La diffusion du match commence à vingt heures.
\end{quote}

\begin{quote}
\emph{The broadcast of the match starts at eight in the evening.}
\end{quote}

\emph{En diffusion} — being shown, on air.

\begin{cousinweb}
\emph{diffusion} is a learned borrowing from Latin \emph{diffūsiō} \textquotedblleft{}a spreading out\textquotedblright{} (EtymDB), from \emph{diffundere} \textquotedblleft{}to pour out\textquotedblright{}.
\end{cousinweb}

\begin{grammarlens}[title={What you've built}]
One more word for telling, explaining and describing at B1.
\end{grammarlens}

\subsection*{Your turn}
\begin{itemize}
\raggedright
  \item \emph{Say it:} \emph{la diffusion}
  \item \emph{Say it:} \emph{la diffusion}, once more
  \item \emph{Say it:} a sentence of your own with \emph{la diffusion}
  \item \emph{Recall:} say the French for short news items, then the French for to broadcast, then say \emph{la diffusion} again
\end{itemize}

\begin{tcolorbox}[breakable,colback=teal!4,colframe=teal!35!black,title={Before you move on}]
What does la diffusion mean? (\textquotedblleft{}The broadcasting, the showing; the spread\textquotedblright{}.) Say it once more.
\end{tcolorbox}
\section[filmer]{filmer — to film}

\label{lesson:FR-C283-filmer}



\begin{quote}
Before the new one: say the French for to broadcast, then the French for the broadcasting.
\end{quote}

\subsection*{You'll want to know: filmer}
\textbf{filmer} — \textquotedblleft{}to film\textquotedblright{}.

\begin{quote}
Un homme a tout filmé avec son téléphone.
\end{quote}

\begin{quote}
\emph{A man filmed everything with his phone.}
\end{quote}

\emph{Se faire filmer} — to get filmed.

\begin{cousinweb}
\emph{filmer} is built on \emph{film}, which was borrowed from English \emph{film} \textquotedblleft{}a thin skin, a coating\textquotedblright{} (EtymDB lists the verb without a parent).
\end{cousinweb}

\begin{grammarlens}[title={What you've built}]
One more word for telling, explaining and describing at B1.
\end{grammarlens}

\subsection*{Your turn}
\begin{itemize}
\raggedright
  \item \emph{Say it:} \emph{filmer}
  \item \emph{Say it:} \emph{filmer}, once more
  \item \emph{Say it:} a sentence of your own with \emph{filmer}
  \item \emph{Recall:} say the French for to broadcast, then the French for the broadcasting, then say \emph{filmer} again
\end{itemize}

\begin{tcolorbox}[breakable,colback=teal!4,colframe=teal!35!black,title={Before you move on}]
What does filmer mean? (\textquotedblleft{}To film\textquotedblright{}.) Say it once more.
\end{tcolorbox}
\section[interviewer]{interviewer — to interview}

\label{lesson:FR-C283-interviewer}



\begin{quote}
Before the new one: say the French for the broadcasting, then the French for to film.
\end{quote}

\subsection*{You'll want to know: interviewer}
\textbf{interviewer} — \textquotedblleft{}to interview\textquotedblright{}.

\begin{quote}
La journaliste a interviewé le maire.
\end{quote}

\begin{quote}
\emph{The journalist interviewed the mayor.}
\end{quote}

It ends in the French \emph{-er} sound of \emph{parler}. The noun is \emph{l'interview}, feminine; in Québec, \emph{l'entrevue}.

\begin{cousinweb}
\emph{interviewer} was borrowed from English \emph{interview}, which itself comes from French \emph{entrevue} \textquotedblleft{}a meeting face to face\textquotedblright{} (EtymDB), built on \emph{voir}, from Latin \emph{vidēre}.
\end{cousinweb}

\begin{grammarlens}[title={What you've built}]
One more word for telling, explaining and describing at B1.
\end{grammarlens}

\subsection*{Your turn}
\begin{itemize}
\raggedright
  \item \emph{Say it:} \emph{interviewer}
  \item \emph{Say it:} \emph{interviewer}, once more
  \item \emph{Say it:} a sentence of your own with \emph{interviewer}
  \item \emph{Recall:} say the French for the broadcasting, then the French for to film, then say \emph{interviewer} again
\end{itemize}

\begin{tcolorbox}[breakable,colback=teal!4,colframe=teal!35!black,title={Before you move on}]
What does interviewer mean? (\textquotedblleft{}To interview\textquotedblright{}.) Say it once more.
\end{tcolorbox}
\section[l'antenne]{l'antenne — the aerial, the antenna; being on air}

\label{lesson:FR-C283-antenne}



\begin{quote}
Before the new one: say the French for to film, then the French for to interview.
\end{quote}

\subsection*{You'll want to know: l'antenne}
\textbf{l'antenne} — \textquotedblleft{}the aerial, the antenna; being on air\textquotedblright{}.

\begin{quote}
Le chanteur a parlé pendant dix minutes à l'antenne.
\end{quote}

\begin{quote}
\emph{The singer spoke on air for ten minutes.}
\end{quote}

\emph{Être à l'antenne} — to be on air. Also an insect's feeler.

\begin{cousinweb}
\emph{antenne} is a learned borrowing from Latin \emph{antemna} \textquotedblleft{}a sail yard\textquotedblright{} (EtymDB), first said of an insect's feelers, then of a radio mast.
\end{cousinweb}

\begin{grammarlens}[title={What you've built}]
One more word for telling, explaining and describing at B1.
\end{grammarlens}

\subsection*{Your turn}
\begin{itemize}
\raggedright
  \item \emph{Say it:} \emph{l'antenne}
  \item \emph{Say it:} \emph{l'antenne}, once more
  \item \emph{Say it:} a sentence of your own with \emph{l'antenne}
  \item \emph{Recall:} say the French for to film, then the French for to interview, then say \emph{l'antenne} again
\end{itemize}

\begin{tcolorbox}[breakable,colback=teal!4,colframe=teal!35!black,title={Before you move on}]
What does l'antenne mean? (\textquotedblleft{}The aerial, the antenna; being on air\textquotedblright{}.) Say it once more.
\end{tcolorbox}
\section[(dialogue)]{Retour sur les chapitres 276–283}

\label{lesson:FR-R283-revision-276-283}



\begin{quote}
Say the French for the time and the decade.
\end{quote}

\subsection*{Your turn}
Say these aloud:

\begin{itemize}
\raggedright
  \item the time, the decade, the millennium, the era, the length
  \item the step, the phase, the chronology, the instant, at the end of
  \item the news, the press, the daily paper, the weekly, the subscriber
  \item to write, the newsroom, to print, to come out, the publication
  \item the television news, the presenter, the viewer, the listener, the microphone
  \item the report, the special correspondent, the correspondent, the dispatch, the news agency
  \item the headline, the banner headline, the section, the editorial, short news items
  \item to broadcast, the broadcasting, to film, to interview, the aerial
\end{itemize}

\begin{tcolorbox}[breakable,colback=teal!4,colframe=teal!35!black,title={Before you move on}]
The time? (\textbf{l'époque}.) The decade? (\textbf{la décennie}.) The millennium? (\textbf{le millénaire}.) The era? (\textbf{l'ère}.) The length? (\textbf{la durée}.) The step? (\textbf{l'étape}.) The phase? (\textbf{la phase}.) The chronology? (\textbf{la chronologie}.) The instant? (\textbf{l'instant}.) At the end of? (\textbf{à l'issue de}.) The news? (\textbf{l'actualité}.) The press? (\textbf{la presse}.) The daily paper? (\textbf{le quotidien}.) The weekly? (\textbf{l'hebdomadaire}.) The subscriber? (\textbf{l'abonné}.) To write? (\textbf{rédiger}.) The newsroom? (\textbf{la rédaction}.) To print? (\textbf{imprimer}.) To come out? (\textbf{paraître}.) The publication? (\textbf{la parution}.) The television news? (\textbf{le journal télévisé}.) The presenter? (\textbf{le présentateur}.) The viewer? (\textbf{le téléspectateur}.) The listener? (\textbf{l'auditeur}.) The microphone? (\textbf{le micro}.) The report? (\textbf{le reportage}.) The special correspondent? (\textbf{l'envoyé spécial}.) The correspondent? (\textbf{le correspondant}.) The dispatch? (\textbf{la dépêche}.) The news agency? (\textbf{l'agence de presse}.) The headline? (\textbf{le gros titre}.) The banner headline? (\textbf{la manchette}.) The section? (\textbf{la rubrique}.) The editorial? (\textbf{l'éditorial}.) Short news items? (\textbf{les faits divers}.) To broadcast? (\textbf{diffuser}.) The broadcasting? (\textbf{la diffusion}.) To film? (\textbf{filmer}.) To interview? (\textbf{interviewer}.) The aerial? (\textbf{l'antenne}.)
\end{tcolorbox}
